% ==============================================================
\section{Homothetic Structure and Wealth Normalization}
\label{sec:dp-homothetic}
% ==============================================================

CES preferences are homothetic: scaling all consumption levels by a constant scales utility by the same constant. In the dynamic problem, this property has a powerful consequence---the value function is linear in wealth, the optimal policy is a constant fraction of wealth, and the state space collapses. The per-period problem becomes a static CES problem of exactly the kind studied in \cref{ch:canonical-problems}, complete with its KL representation. This section develops the implication chain.

% --------------------------------------------------------------
\subsection{Homotheticity of CES Preferences}
\label{subsec:dp-homotheticity}
% --------------------------------------------------------------

Recall that the recursive utility $\V_t = \Mbar(c_t, \V_{t+1}; 1-\beta, \rho)$ inherits the homogeneity of $\Mbar$.

\paragraph{Homogeneity of the Aggregator.}
The classical mean $\Mbar(x_1, x_2; \omega, \rho) = (\omega x_1^\rho + (1-\omega) x_2^\rho)^{1/\rho}$ is homogeneous of degree one:
\begin{equation}
\Mbar(\lambda x_1, \lambda x_2; \omega, \rho) = \lambda \, \Mbar(x_1, x_2; \omega, \rho), \quad \lambda > 0.
\label{eq:mbar-homogeneous}
\end{equation}
Iterating, the recursive utility $\V_0$ is homogeneous of degree one in the entire consumption stream: scaling $\{c_t\}$ by $\lambda$ scales $\V_0$ by $\lambda$.

\paragraph{Homogeneity of the Constraint Set.}
The transition $W' = R(W - c)$ is also linear, so the feasible set of consumption streams scales linearly with $W_0$: from initial wealth $\lambda W_0$, the agent can attain exactly $\lambda$ times any stream attainable from $W_0$. Combined with homogeneity of preferences, this delivers the following.

\begin{proposition}[Homogeneity of the Value Function]
\label{prop:V-homogeneous}
The optimal value function $V$ in the deterministic problem (and in the stochastic problem with i.i.d.\ returns or, more generally, returns whose distribution is independent of $W$) satisfies
\begin{equation}
V(\lambda W) = \lambda V(W), \quad \lambda > 0.
\label{eq:V-homogeneous}
\end{equation}
\end{proposition}

\begin{proof}[Proof sketch]
By induction on the finite-horizon problem and passage to the limit, or directly via the principle of optimality: if $\{c_t^\star\}$ is optimal from $W$, then $\{\lambda c_t^\star\}$ is feasible from $\lambda W$, and homogeneity of preferences gives a value of $\lambda V(W)$. The optimality of $\{\lambda c_t^\star\}$ from $\lambda W$ follows because any improvement could be scaled down to improve $\{c_t^\star\}$, contradicting its optimality.
\end{proof}

% --------------------------------------------------------------
\subsection{The Normalized Bellman Equation}
\label{subsec:dp-normalized-bellman}
% --------------------------------------------------------------

Homogeneity allows us to write $V(W) = v \cdot W$ for some scalar $v$, eliminating wealth from the Bellman equation.

\paragraph{Substitution.}
Using $V(W) = v W$ and $V(W') = v W'$ in the Bellman equation \cref{eq:bellman-equation}:
\begin{equation}
v W = \max_{c \in [0, W]} \Mbar\Big( c, \, \mu(v R'(W - c)); \; 1-\beta, \, 1-1/\theta \Big).
\label{eq:bellman-vW}
\end{equation}

Pull out the random gross return: $\mu(v R'(W-c)) = (W-c) \, v \, \mu(R')$ if $v$ and $W-c$ are deterministic and $\mu$ is positively homogeneous. (The certainty equivalent $\mu(X) = (\E[X^{1-\gamma}])^{1/(1-\gamma)}$ is indeed homogeneous of degree one.) Then:
\begin{equation}
v W = \max_{c} \Mbar\Big( c, \, (W-c) \, v \, \mu(R'); \; 1-\beta, \, 1-1/\theta \Big).
\label{eq:bellman-pulled}
\end{equation}

\paragraph{Re-Parameterizing.}
Let $m = c/W$ denote the consumption-to-wealth ratio. Then $c = mW$ and $W - c = (1-m)W$. Substituting and using homogeneity of $\Mbar$:
\begin{equation}
v W = W \cdot \max_{m \in [0, 1]} \Mbar\Big( m, \, (1-m) \, v \, \mu(R'); \; 1-\beta, \, 1-1/\theta \Big).
\label{eq:bellman-m}
\end{equation}

Dividing by $W$ delivers the \emph{normalized Bellman equation}:
\begin{equation}
\boxed{v = \max_{m \in [0, 1]} \Mbar\Big( m, \, (1-m) \, v \, \mu(R'); \; 1-\beta, \, 1-1/\theta \Big).}
\label{eq:normalized-bellman}
\end{equation}
A scalar fixed-point equation in the scalar unknown $v$. Wealth has disappeared.

\begin{calloutimportant}[The State-Space Collapse]
The state variable $W$ has been eliminated. The Bellman equation, originally a functional equation for $V(\cdot)$ on $\mathbb{R}_+$, has become a single algebraic equation in a single number $v$. This is the practical payoff of CES homotheticity: a problem that would generically be infinite-dimensional is solved by characterizing one scalar.
\end{calloutimportant}

\paragraph{Optimal Consumption Rate.}
The maximizer of \cref{eq:normalized-bellman}, denoted $m^\star$, is the optimal consumption-to-wealth ratio. The first-order condition is:
\begin{equation}
(1-\beta) \, m^{-1/\theta} = \beta \, v \, \mu(R') \cdot \big( (1-m) v \mu(R') \big)^{-1/\theta}.
\label{eq:m-foc}
\end{equation}
Rearranging:
\begin{equation}
\frac{m^\star}{1 - m^\star} = \left( \frac{1-\beta}{\beta} \right)^\theta \cdot \big( v \, \mu(R') \big)^{\theta - 1}.
\label{eq:m-formula}
\end{equation}
Substituting back into \cref{eq:normalized-bellman} pins down $v$.

\paragraph{Closed Form.}
After algebra, the stationary normalized value satisfies
\begin{equation}
v = \left( (1-\beta)^\theta + \beta^\theta \mu(R')^{\theta - 1} \right)^{1/(\theta - 1)} \cdot (1-\beta)^{\theta/(1-\theta)},
\label{eq:v-closed}
\end{equation}
and the consumption rate is
\begin{equation}
m^\star = 1 - \beta^\theta \mu(R')^{\theta - 1} (1-\beta)^{1-\theta}.
\label{eq:m-closed}
\end{equation}
For $\theta = 1$ (log utility), \cref{eq:m-closed} simplifies to the classic result $m^\star = 1 - \beta$: the agent consumes a constant fraction $1-\beta$ of wealth each period, regardless of the return.

% --------------------------------------------------------------
\subsection{Recovering the Static Canonical Problem}
\label{subsec:dp-static-recovery}
% --------------------------------------------------------------

The normalized Bellman equation \cref{eq:normalized-bellman} is, structurally, the static two-good CES problem of \cref{ch:canonical-problems}.

\paragraph{The Mapping.}
Compare the canonical static problem
\begin{equation}
U^\star = \max_{x_1 + x_2 = 1} \Mbar(x_1, x_2; \omega, \rho)
\label{eq:static-canonical}
\end{equation}
(after normalizing the budget to one) with the normalized Bellman \cref{eq:normalized-bellman}, which fixes the resource as $1$ (since $m + (1-m) = 1$). The mapping is:
\[
x_1 \leftrightarrow m, \quad x_2 \leftrightarrow 1 - m, \quad \omega \leftrightarrow 1 - \beta,
\]
together with a shadow price $p_1 = 1$ on consumption today and a shadow price $p_2 = 1/(v \mu(R'))$ on consumption tomorrow.

\paragraph{The Shadow Price of the Future.}
The shadow price of next-period ``goods'' is
\begin{equation}
p_2 = \frac{1}{v \, \mu(R')}.
\label{eq:shadow-price}
\end{equation}
This is the cost (in today's units) of one unit of certainty-equivalent normalized future wealth. Three pieces enter: the inverse of the certainty equivalent of returns ($1/\mu(R')$), the inverse of the per-unit value of future wealth ($1/v$), and the implicit normalization $p_1 = 1$.

\begin{calloutnote}[The Dynamic Price Index]
The shadow price $p_2$ plays the role of the price index $\Pstar$ from \cref{ch:comparative-statics}, but for the dynamic problem. It aggregates the random return and the per-unit continuation value into a single number that prices ``the future.'' The KL representation of welfare developed in \cref{sec:price-index-kl} carries over directly: the agent's value $v$ depends on the dispersion of $\mu(R')$ across states, with the dispersion penalty governed by KL divergence.
\end{calloutnote}

% --------------------------------------------------------------
\subsection{The KL Representation in Dynamic Settings}
\label{subsec:dp-kl-representation}
% --------------------------------------------------------------

The KL representation of CES welfare (\cref{ch:information}) extends to dynamic problems with no modification beyond the relabeling above.

\paragraph{Recap: The Static Result.}
For a static CES problem with prices $\bp$ and weights $\bom$, the welfare loss from price dispersion is governed by the KL divergence between the optimal expenditure-share weights $\bw^\star$ and the preference weights $\bom$:
\begin{equation}
\frac{d \log \Pstar}{d \tilde{\rho}} = \frac{\KL{\bw^\star}{\bom}}{\tilde{\rho}^2}, \qquad \tilde{\rho} = \frac{\rho}{\rho - 1}.
\label{eq:kl-static}
\end{equation}
Higher KL divergence means greater welfare sensitivity to substitutability.

\paragraph{Two-Period Dynamic Case.}
For the two-period problem (today vs.\ tomorrow), the KL divergence is between the weights on $c_t$ and $c_{t+1}$ implied by the budget at the optimum, and the preference weights $(1-\beta, \beta)$. With deterministic return $R$:
\begin{equation}
w_1^\star = \frac{c_t}{c_t + c_{t+1}/R}, \quad w_2^\star = 1 - w_1^\star, \quad \omega_1 = 1-\beta, \quad \omega_2 = \beta.
\label{eq:dynamic-kl-weights}
\end{equation}
The KL divergence $\KL{\bw^\star}{\bom}$ measures the gap between optimal allocation shares and preference shares; it vanishes when $\beta R = 1$ (Euler equation balanced) and grows otherwise.

\paragraph{Stochastic Case.}
With random returns, the KL representation includes an additional layer: the divergence between the risk-neutral weights induced by the certainty equivalent and the preference weights. The decomposition
\begin{equation}
\frac{d \log v}{d \tilde\rho} = \KL{\text{intertemporal}}{\text{preferences}} + \KL{\text{across states}}{\text{probabilities}}
\label{eq:kl-decomposition-dyn}
\end{equation}
separates the welfare cost of price dispersion across time from the welfare cost of return dispersion across states. The first term governs the IES; the second governs risk aversion. They are decoupled in Epstein--Zin precisely because the inner and outer aggregators have different parameters.

\begin{calloutimportant}[Unification]
Three observations come together at this point:
\begin{enumerate}[leftmargin=2em]
    \item The dynamic problem is solved by a normalized Bellman equation with no state variable.
    \item The normalized Bellman equation is, structurally, a static two-good CES problem.
    \item All comparative statics, KL representations, and welfare formulas from the static chapter apply with appropriate relabeling.
\end{enumerate}
The dynamic CES problem is, in a precise sense, the static CES problem repeated. This is the theoretical basis for the applications developed in subsequent chapters---consumption and savings, portfolio choice, and labor supply---each of which is, internally, a static CES allocation.
\end{calloutimportant}

% --------------------------------------------------------------
\subsection{When Homotheticity Fails}
\label{subsec:dp-non-homothetic}
% --------------------------------------------------------------

Several extensions break homotheticity and require the full state-space treatment.

\paragraph{Borrowing Constraints.}
A constraint $W - c \geq \underline{a}$ for some absolute floor $\underline{a} > 0$ violates the linear scaling of the constraint set. Wealth re-enters as a state variable, the consumption rate $m$ varies with $W$, and the value function is no longer linear.

\paragraph{Subsistence Consumption.}
A specification $u(c - \bar{c})$ for some subsistence level $\bar{c} > 0$ breaks homogeneity at low wealth levels. The agent's effective discount rate and IES vary with wealth.

\paragraph{Labor Income.}
With non-zero labor income $y_t$, total resources are $W_t + y_t/R + \cdots$, and the relevant ``wealth'' is human plus financial. If labor income is bounded away from zero, the same homothetic argument applies to total wealth, but only in the long-horizon limit.

\paragraph{Time-Varying Investment Opportunities.}
If the distribution of returns depends on a state $z_t$ (e.g., a stochastic volatility process), the value function becomes $V(W, z) = v(z) \cdot W$ with $v$ a non-trivial function of $z$. Homotheticity in $W$ survives, but a separate fixed-point problem in $v(z)$ remains.

\begin{calloutnote}[Recovering Tractability]
In each of these cases, the KL representation and the canonical-problem mapping continue to hold \emph{conditional on the additional state}. The role of CES is to provide a one-dimensional reduction in $W$, leaving the additional state(s) as the only sources of complexity. This is the divide-and-conquer logic that makes CES-based dynamic models tractable even when full homotheticity fails.
\end{calloutnote}

% --------------------------------------------------------------
\subsection{Summary}
\label{subsec:dp-homothetic-summary}
% --------------------------------------------------------------

\begin{calloutimportant}[Key Results]
\begin{enumerate}[label=(\roman*), leftmargin=2em]
    \item CES preferences and a linear budget constraint imply a \textbf{homogeneous value function}: $V(W) = v \cdot W$ for a scalar $v$.
    
    \item The Bellman equation reduces to a \textbf{normalized Bellman equation} in the consumption rate $m$ and the per-unit value $v$, with $W$ eliminated.
    
    \item Closed-form solutions: $m^\star = 1 - \beta^\theta \mu(R')^{\theta-1} (1-\beta)^{1-\theta}$ and a corresponding closed form for $v$.
    
    \item The normalized problem has \textbf{the same structure as a static two-good CES problem}, with the gross return playing the role of an inverse price.
    
    \item The \textbf{KL representation} of welfare from \cref{ch:information} carries over to the dynamic problem, with separate KL terms for intertemporal and cross-state dispersion in the Epstein--Zin case.
    
    \item Homotheticity fails under borrowing constraints, subsistence consumption, or non-scalar wealth, but the canonical-problem mapping survives \emph{conditionally} on additional state variables.
\end{enumerate}
\end{calloutimportant}
